\subsection{Cotas uniformes de los pesos y de las vacancias; interioridad}\label{sec:P02}
Se mantienen las notaciones de la Sección~\ref{sec:P01}, en particular
$n\ge32$, $\kappa=1/20$, $M_0=6/5$, $d_0=3\log(4/3)$ y
$T=(M_0+d_0+4/e)/\kappa$.

\paragraph{Factorización con el multiplicador presente.}
La ecuación (P1.3) admite la factorización
\[
 x_{ij}=B R_i C_j K_{ij},\qquad
 K_{ij}=h_{D(i,j)}h_{A(i,j)}e^{-\mu_{ij}}(1-p_{ij})^{-3},
                                                               \tag{P2.1}
\]
con factores positivos y normalizaciones geométricas
\[
 \frac1n\sum_i\log R_i=\frac1n\sum_j\log C_j=0.
\]
Se obtienen, por ejemplo, dividiendo $e^{-\alpha_i}$ y
$e^{-\beta_j}$ por sus medias geométricas y absorbiendo en $B$
esos dos factores escalares junto con $ne^{-2}$.
Cada fila encuentra como máximo una vez cada diagonal de cada
familia. Por tanto, (P1.5) y la desigualdad $(1-p)^{-3}\ge1$ dan
\[
 \frac1n\sum_j\log K_{ij}\ge-T.
\]
La misma propiedad vale para cada columna:
\[
 \frac1n\sum_i\log K_{ij}\ge-T.                  \tag{P2.2}
\]
Aquí se usa la suma global de $\mu$: la suma de sus valores en una
fila o columna es menor o igual a esa suma global.

\paragraph{Control de los factores sin suponer vacancia puntual.}
La desigualdad de Jensen, las normalizaciones geométricas y las sumas de filas dan
\[
 1=B R_i\frac1n\sum_j C_jK_{ij}\ge B R_i e^{-T}.
\]
Aplicando lo mismo a las columnas,
\[
 R_i,C_j\le e^T/B.                               \tag{P2.3}
\]
Por otro lado $\log K_{ij}\le d_0$, pues $h_d\le1$,
$\mu_{ij}\ge0$ y $p_{ij}\le1/4$.
Multiplicamos el logaritmo de (P2.1) por $p_{ij}$ y sumamos.
Las contribuciones de $\log R_i$ y $\log C_j$ desaparecen: las
sumas de filas y columnas son uno y sus sumas de logaritmos son
cero. Por la convexidad ya usada en la Sección~\ref{sec:P01},
\[
 n\log B=\sum_{ij}p_{ij}\log(np_{ij})
               -\sum_{ij}p_{ij}\log K_{ij}\ge-d_0n.
\]
Así $B\ge e^{-d_0}$. De (P2.3) se deduce la cota
\[
 B R_i C_j(1-p_{ij})^{-3}
 \le e^{2T+d_0}/B\le e^{2T+2d_0}=:M.             \tag{P2.4}
\]
En (P2.1) todos los factores $h_d,e^{-\mu_{ij}}$ son a lo sumo
uno. Por ello, todavía sin retirar los topes,
\[
             0<p_{ij}\le M/n.                    \tag{P2.5}
\]

\paragraph{Vacancias uniformes.}
Fijada una diagonal $d$ de la primera familia, (P2.1) y (P2.4) dan
\[
 m_d(p)=\frac{h_d}{n}\sum_{(i,j)\in d}
 B R_i C_j h_{A(i,j)}e^{-\mu_{ij}}(1-p_{ij})^{-3}
 \le M h_d\frac{|d|}{n}\le M h_d.
\]
El mismo argumento se aplica a la segunda familia.
Como $m_d=1-h_d$, resulta para todas las diagonales
\[
                 h_d\ge\eta:=\frac1{1+M}.        \tag{P2.6}
\]
También esta conclusión es anterior a la retirada del tope.

\paragraph{Retirada del tope y cota inferior.}
Para
\[
 n\ge n_0:=\max\{32,\lfloor4M\rfloor+1\},
\]
(P2.5) implica $p_{ij}<1/4$. La complementariedad da entonces
$\mu_{ij}=0$ en todas las casillas.
Así se obtienen la hipótesis $(H_\infty)$ de la Definición~\ref{def:Hinf} y
las cotas uniformes. De hecho, las cotas (P2.5) y
(P2.6) valen ya para todo $n\ge32$.
Para $n\ge n_0$ se tiene además una cota inferior de orden $1/n$.
De $K_{ij}\le e^{d_0}$ y (P2.3), la ecuación de fila implica
\[
 R_i=\left(B\frac1n\sum_jC_jK_{ij}\right)^{-1}
       \ge e^{-(T+d_0)};
\]
análogamente $C_j\ge e^{-(T+d_0)}$.
Una vez que $\mu=0$, (P2.6) da $K_{ij}\ge\eta^2$. Junto con
$B\ge e^{-d_0}$, esto prueba
\[
 \frac{c}{n}\le p_{ij}\le\frac Mn,\qquad
 c:=\eta^2e^{-2T-3d_0}>0,
 \qquad h_d\ge\eta.                              \tag{P2.7}
\]
